\documentclass[english, parskip=full]{scrartcl}
\usepackage[utf8]{inputenc}  % Kodierung der Datei
\usepackage[english]{babel}  % Sprache auf Deutsch 
\usepackage{color}
\usepackage{graphicx, gensymb}
\usepackage{amsfonts, cancel, amssymb, slashed}
\usepackage{amsmath, amsthm, array, esvect, esint}
\usepackage{booktabs, multirow}  % schönere Tabellen
\usepackage{caption}  % Anpassung der Captions
\usepackage{scrlayer-scrpage}    % Manipulation des Seitenstils
\pagestyle{scrheadings}  % Stil für die Seite setzen
\automark{section}  % Abschnittsnamen als Seitenbeschriftung 
\setheadsepline{.5pt}    % Kopzeile durch Linie abtrennen
\usepackage[hidelinks]{hyperref}  % Links und weitere PDF-Features
\usepackage{typearea, tikz, pgffor, verbatim}
\usetikzlibrary{calc, quotes}
\usepackage[cache=false, outputdir=build]{minted}
\usepackage{dsfont}


\definecolor{bg}{rgb}{0.9,0.9,0.9}
\newcommand\numberthis{\addtocounter{equation}{1}\tag{\theequation}}
\newcommand{\bra}{}
\newcommand{\kom}{}
\newcommand{\brak}{}
\newcommand{\braket}{}
\newcommand{\intx}{}
\newcommand{\intr}{}
\newcommand{\kla}{}
\newcommand{\klam}{}
\newcommand{\klammer}{}
\newcommand{\oper}{}
\newcommand{\tecxt}{}
\newcommand{\Bra}{}
\newcommand{\Ket}{}
\def\I{\text{i}}
\def\E{\text{e}}

\newcommand*{\titel}{Formulas for Impress Presentation}
\newcommand*{\autor}{Joachim Schwardt}

\hypersetup{pdfauthor={\autor}, pdftitle={\titel}}

\subject{Bachelor thesis}
\title{\titel}
\author{\autor}
\date{\today}

\begin{document}

%\begin{titlepage}
\maketitle  % Titel setzen
%\tableofcontents  % Inhaltsverzeichnis setzen
%\end{titlepage}

\section{Frequency analysis}
\begin{align*}
w_{n,N}&=\frac{1}{S}g\kla\left( {\frac{n}{N}} \right)&\tecxt\text{with}&&S&=\sum_{n=0}^{N-1}g\kla\left( {\frac{n}{N}} \right),
\end{align*} 
using an infinitely smooth function $g\in C^\infty_c (\mathbb{R})$ with compact support
\begin{align*}
g(n)&=\begin{cases}
\exp\kla\left( {-\klam\left[ {n(1-n)} \right]^{-1}} \right)&\tecxt\text{for}\ 0<n<1,\\
0&\tecxt\text{else}
\end{cases}
\end{align*}
By construction, these weights are normalized, i.e. $$\sum_{n=0}^{N-1}w_{n,N}=1$$
Then the weighted Birkhoff average can be defined as
\begin{align*}
\tecxt\text{WA}_N(X_N)&=\sum_{n=0}^{N-1}w_{n,N}x_n
\end{align*}
Naff
\begin{align*}
z_n&\approx\sum_{k= 1}^m a_k\E^{2\pi\I\nu_k n}
\end{align*}
\begin{align*}
\phi(\nu)&=\sum_{n=0}^{N-1}\chi_n z_n \E^{-2\pi\I\nu n}
\end{align*}
That is, $\nu_1=\tecxt\text{argmax\,}\vert\phi(\nu)\vert$.
Here the 
$$\chi_n=\frac{2}{N}\sin^2\kla\left( {\frac{\pi n}{N}} \right)$$ 
are normalized weights, similar to the $w_n$ used for the WBA.
\section{Two dimensions}
$$f_\nu(q)=q+\nu\,\tecxt\text{mod}\, 1$$ 
for 
$$q\in [0,1)$$ 
with a frequency 
$$\nu\in [0,1)$$
Input coordinate for WBA
$$x_n=(q_{n+1}-q_n)\, \tecxt\text{mod}\, 1$$
Oscillatory signal for $n=0,\dots,N-1$ and $R>0$ and $A<1$
\begin{align*}
z_n&=R\E^{2\pi\I\nu n}=q_n+\I p_n\\
z_n&=R\E^{2\pi\I\nu_1 n}+A\cdot R\E^{2\pi\I\nu_2 n}
\end{align*}
Complex signal
\begin{align*}
z_n&=q_n+\I p_n
\end{align*}
Polar angle
\begin{align*}
\phi_n&=\frac{1}{2\pi}\tecxt\text{arctan2}(q_n,p_n)
\end{align*} 
Angle differences
\begin{align*}
\Delta\phi_n&=\phi_{n+1}-\phi_n
\end{align*} 
and new input coordinate for WBA 
$$x_n=\Delta\phi_n\,\tecxt\text{mod}\,1$$
Number of correct digits
\begin{align*}
\tecxt\text{dig}_N&=-\log_{10}|\nu-\nu_N|
\end{align*}
Chirikov´s two dimensional standard map
\begin{align*} 
q_{n+1}&=q_{n}+p_{n}\\
p_{n+1}&=p_{n}+\frac{K}{2\pi}\sin (2\pi q_{n+1})
\end{align*}
Here $K=0.9$.
Chaos indicator
\begin{align*}
\Delta\nu (N) &=|\nu_{[0,N-1]}-\nu_{[N,2N-1]}|
\end{align*} 
Time-dependent frequency of trapped orbit
\begin{align*}
(\nu_n)&=\klammer\left\{\nu_{[n\sigma,n\sigma + L-1]}\,\middle\vert\,n=0,1,2,\dots,\left\lfloor\frac{N-L}{\sigma}\right\rfloor \right\}
\end{align*}
\section{Four dimensions}
Four dimensional standard map
\begin{align*}
q_{1,n+1}&=q_{1,n}+p_{1,n}\\
p_{1,n+1}&=p_{1,n}+\frac{K_1}{2\pi}\sin\kla\left( {2\pi q_{1,n+1}} \right)+C\\
q_{2,n+1}&=q_{2,n}+p_{2,n}\\
p_{2,n+1}&=p_{2,n}+\frac{K_2}{2\pi}\sin\kla\left( {2\pi q_{2,n+1}} \right)+C\\
C&=\frac{K}{2\pi}\sin\kla\left( {2\pi[q_{1,n+1}+q_{2,n+1}]} \right)
\end{align*}
Set of points
$$X_N=\{\mathbf{x}_n\in\mathbb{R}^d\,\vert\,n=0,\dots,N-1\}$$
Inertia tensor
\begin{align*}
I_d(X_N)&=\sum_{n=0}^{N-1}\mathbf{x}_n\cdot \mathbf{x}_n \mathds{1}_{d\times d}-\mathbf{x}_n\otimes \mathbf{x}_n. \numberthis\label{eqn:InertiaTensor}
\end{align*}
Orthogonal transformation matrix 
$$U\in \text{O}(d)$$
satisfying 
$$I_d=U DU^\top$$
Transformed points
\begin{align*}
\mathbf{x}'_n&=U\mathbf{x}_n
\end{align*}
Ratio test
\begin{align*}
r_{jk}&=\frac{\tecxt\text{min}\klammer\left\{ \sqrt{x_{j,n}^2+x_{k,n}^2} \,\middle\vert\,n=0,\dots,N-1 \right\}}{\tecxt\text{max}\klammer\left\{\sqrt{x_{j,n}^2+x_{k,n}^2} \,\middle\vert\,n=0,\dots,N-1 \right\}}
\end{align*}
Resonance condition for integers $m_1,m_2,l\in\mathbb{Z}$
\begin{align*}
m_1\nu_1+m_2\nu_2&=l
\end{align*}
Uni-modular transformation
\begin{align*}
\begin{pmatrix}
\nu_1 \\ \nu_2
\end{pmatrix}&\longrightarrow\begin{pmatrix}
1&0 \\ 4& -1
\end{pmatrix}\begin{pmatrix}
\nu_1 \\ \nu_2
\end{pmatrix}, \numberthis\label{eqn:UnimodularTransformation}
\end{align*}
for all black points $(\nu_1,\nu_2)$ on the $3:1:1$ resonance.\\
Frequency space $10^8$ random initial conditions with $q\in [0.2, 0.8]$ and $p\in [-0.3, 0.3]$.
$N=2048$ or $N=4096$ with $K_1=2.25, K_2=3$ and $K=1$.
\section{Appendix}
parameterized transformation
\begin{align*}
\Delta\phi_n&\longrightarrow\klam\left[ {\kla\left( {\Delta\phi_n+\alpha} \right)\,\tecxt\text{mod}\,1} \right]-\alpha
\end{align*}
WBA for dual frequency signals
Angle differences
\begin{align*}
\hat{\Delta}\phi_n&=\Delta\phi_n -k_n
\end{align*}
for integers $k_n\in\mathbb{Z}$ such that
$\hat{\Delta}\phi_n$ lie within an interval of length one.
Visual splitting now described by $k_n\neq k_m$ for some $n,m$.
\\
Angle difference $\Delta\phi_m$ as point of reference, i.e. choose $k_m=0$.
Differences
\begin{align*}
\delta_{nm}&=\Delta\phi_n - \Delta\phi_m&\tecxt\text{and}&&\hat{\delta}_{nm}&=\hat{\Delta}\phi_n - \hat{\Delta}\phi_m
\end{align*}
related by
\begin{align*}
\delta_{nm}&=\hat{\delta}_{nm}+k_n
\end{align*}
Compute $k_n$ based on rules
\begin{align*}
k_n&=\begin{cases}
+1&\tecxt\text{if}\ \delta_{nm}>1/2 \\
-1&\tecxt\text{if}\ \delta_{nm}<-1/2 \\
\ 0 &\tecxt\text{else}
\end{cases}
\end{align*}
\end{document}